% !TEX root = ../nc-quadric-moduli.tex
% 2026-11-21 / v1.39
\section{noncubic 境界の適合根と Hochschild 作用}
\label{sec:noncubic-hh-roots}


\paragraph{結論と範囲.}
中央二重平方表示の二つの根対を交換する\emph{trace 零の}線型対合を選ぶと、
対応する適合 Serre 根は有限 foundation の $\operatorname{HH}^1,\operatorname{HH}^2$ に恒等に作用する。
ただし根対が一致するときは対角の交換対合を選ぶ。
trace 零だけでは十分でないことと、選択条件の具体的な行列式を後で示す。
この計算は noncubic 表示の一次元の変形だけを扱うものではない。
foundation の全変形接空間、すなわち根対が異なる場合の三次元、
一致する場合の四次元を扱う。
従って既存の Artin 根持上げ補題を適用すると、指定した良い根は
全ての Artin Abel 変形へ持ち上がり、その Morita 類は一意である。
さらに有限 foundation の各 noncubic 点には、
この根と実際の Serre 平方同型を持つ有限型 étale 近傍が存在する。
同じ根に対する AS4 被覆の横断変形への持上げと、有限型基底上の全相対認識は
本節の結論に含めない。
本節の計算の新規性は未確認である。

\subsection{四項 foundation の二項移動}

$k$ を標数零の代数閉体とする。
full geometric helix の四項を $E_0,E_1,E_2,E_3$ とし、
\[
A=\Hom(E_0,E_1),\quad B=\Hom(E_1,E_2),\quad
C=\Hom(E_2,E_3),\quad D=\Hom(E_0,E_2)
\]
と置く。$A,B,C$ は二次元、$D$ は四次元である。
$W=\Hom(E_0,E_3)$ は六次元であり、
\[
F:A\otimes B\longrightarrow D,\qquad
G:C\otimes D\longrightarrow W
\]
が foundation の積を与える。
因子の並べ替えを省略して書く。
$G$ は全射とし、$K=\ker G$ は二次元とする。
包含 $K\subset C\otimes D$ を curry した写像を
\[
j:D^{\vee}\longrightarrow C\otimes K^{\vee}
\]
と書く。さらに
\begin{equation}\label{eq:nchh-delta}
\delta=(-F^{\vee},j):D^{\vee}\longrightarrow
H=(A^{\vee}\otimes B^{\vee})\oplus(C\otimes K^{\vee})
\end{equation}
は単射とする。これらは、以下で使う noncubic 点の近傍で開条件である。
$\mathscr L=\operatorname{im}\delta\in\Gr(4,H)$ と置く。

\begin{lemma}[二項移動の有限線型代数]\label{lem:nchh-gale}
$E_{i+4}=S^{-1}E_i[2]$ により延長し、四項窓を
$(E_2,E_3,E_4,E_5)$ へ移す。
対応するデータは
\begin{align}
A'&=C,& B'&=K^{\vee},& C'&=A,\label{eq:nchh-new-small}\\
D'&=H/\mathscr L,&
F'&:C\otimes K^{\vee}\longrightarrow D',\label{eq:nchh-new-D}\\
K'&=\operatorname{im}\bigl(B^{\vee}\longrightarrow A\otimes D'\bigr),\label{eq:nchh-new-K}
\end{align}
である。\eqref{eq:nchh-new-K} は $A$ の coevaluation と第一 summand の商写像の合成である。
$K'\simeq B^{\vee}$ となる開近傍で、この新しい四平面は
\begin{equation}\label{eq:nchh-gale}
\mathscr L'=\{(-Y,X):(X,Y)\in\mathscr L^{\perp}\}
 \subset(C^{\vee}\otimes K)\oplus(A\otimes B)
\end{equation}
である。ここで $\perp$ は $H$ と $H^{\vee}$ の標準評価による annihilator を表す。
この記述は任意の局所 Artin 変形にも同じ式で適用できる。
\end{lemma}

\begin{proof}
符号も含めて相対 bar 複体で確認する。
積を $\Gamma_{ij}=\Hom(E_i,E_j)$ とする上三角代数として書き、
左 projective $P_i=\Gamma e_i$ を使えば $E_i$ に対応する。
$I_i=\Hom_k(e_i\Gamma,k)=S(P_i)$ と置くと
\[
\Hom(E_j,E_4)=\Ext^2_{\Gamma}(I_j,P_0).
\]
$I_3$ は頂点三の単純加群であり、bar 複体の可消去な成分を除くと
\[
\Ext^2(I_3,P_0)=\operatorname{coker}(G^{\vee})=K^{\vee}
\]
である。

$I_2$ は頂点二で $k$、頂点三で $C^{\vee}$ を持ち、
$C$ の作用は評価である。
bar cochain の次数一は
$D^{\vee}\oplus\Hom(W,C)$、次数二は
\[
(A B)^{\vee}\oplus\Hom(A B C,C)\oplus\Hom(D C,C),
\]
次数三は $\Hom(A B C,C)$ である。
次数二の元を $(r,g,h)$ とすると、次の微分は
\[
(r,g,h)\longmapsto g-h(F\otimes1_C)-r\otimes1_C.
\]
従って kernel では中央成分 $g$ を消去できる。
次数一の $(a,b)$ の像は、残る二成分で
\[
(-aF,-bG+a\otimes1_C).
\]
$bG$ を除くことは $h$ を $K$ へ制限することに等しく、
残る $a$ の像はちょうど $(-F^{\vee}a,j(a))$ である。
従って \eqref{eq:nchh-new-D} を得る。
前合成 $E_2\to E_3$ は第二 summand の包含なので、$F'$ も表示した写像となる。

次に $0\to A\otimes P_0\to P_1\to S_1\to0$ を用いる。
頂点一、二、三の部分代数は道代数なので
$\Ext^1(I_2,S_1)=B^{\vee}$、$\Ext^2(I_2,S_1)=0$ である。
長完全列における境界写像は
\[
B^{\vee}\longrightarrow A\otimes\Ext^2(I_2,P_0)=A\otimes D'
\]
であり、bar 複体で代表を取れば
$b^{\vee}\mapsto\sum_a a\otimes[a^{\vee}\otimes b^{\vee}]$ と計算される。
これが \eqref{eq:nchh-new-K} である。
$G'$ は $A\otimes D'$ をその像で割った六次元商である。
また $\Hom(E_4,E_5)=A$ は Serre 移動による同定を使う。

$D'^{\vee}$ は $\mathscr L^{\perp}\subset
(A\otimes B)\oplus(C^{\vee}\otimes K)$ である。
$F'^{\vee}$ は第二成分 $Y$ への射影、$j'$ は第一成分 $X$ への射影だから、
$(-F'^{\vee},j')$ の像は \eqref{eq:nchh-gale} となる。
全て有限 projective な複体と直和因子の消去であるから、
閉ファイバーで必要な階数が成立する任意の Artin 変形でも同じ計算が成立する。
\end{proof}

\begin{remark}[全変形空間を使うこと]
$G$ の商の基底を忘れ、代わりに $K$ の基底を選ぶと、
foundation のデータは $(F,j)$ である。
$\delta$ が単射である開部分で $D$ の基底変更を除くことは、
四平面 $\mathscr L\subset H$ を取ることに等しい。
従って全 foundation 変形の局所模型は
$\Gr(4,8)$ の開部分を
\[
\widetilde G=\GL(A)\times\GL(B)\times\GL(C)\times\GL(K)
\]
で割ったものとなる。
第一 summand への作用は $g_A^{-t}\otimes g_B^{-t}$、
第二 summand への作用は $g_C\otimes g_K^{-t}$ である。
自明に作用する三次元中心を除いた群を $G_{\mathrm{eff}}$ と書く。
その次元は十三で、Grassmannian の次元は十六である。
この模型には noncubic 表示を保たない全ての横断方向が含まれている。
\end{remark}

\subsection{中央二重平方と良い根の選択}

\begin{lemma}[trace 零の交換対合]\label{lem:nchh-exchange}
二つの非退化二元二次形式 $f_1,f_2$ に対し、座標とそれぞれの平方根変数の
定数倍を選び、
\begin{equation}\label{eq:nchh-normal-forms}
f_1=uv,\qquad f_2=\ell_1\ell_2,
\qquad
\begin{pmatrix}\ell_1\\\ell_2\end{pmatrix}
=L_0\begin{pmatrix}u\\v\end{pmatrix},\qquad
L_0=\begin{pmatrix}\alpha&\beta\\\gamma&-\alpha\end{pmatrix},\quad
\alpha^2+\beta\gamma=1
\end{equation}
とできる。二根対が一致する場合にも $L_0=\operatorname{diag}(1,-1)$ と選べる。
\end{lemma}
\begin{proof}
既存の 補題~\ref{lem:noncubic-local-exchange} の行列
$\left(\begin{smallmatrix}1&-a\\b&-1\end{smallmatrix}\right)$ を使う。
その平方は単元 $(1-ab)I$ であり、平方根で割れば trace 零の対合となる。
$L_0f_1=\eta f_2$ の単元 $\eta$ を $d$ の定数倍に吸収する。
この構成は二根対一致の場合にも $a=b=0$ で成立する。
$f_1$ の非退化性により $f_1=uv$ と正規化できる。
\end{proof}

以下、
\[
\mathcal B=k[u,v]\langle c,d\rangle/(c^2-uv,d^2-\ell_1\ell_2)
\]
と置く。$u,v$ は中央である。
\begin{equation}\label{eq:nchh-sigma}
\sigma(u,v)=L_0(u,v),\qquad\sigma(c)=d,\qquad\sigma(d)=-c
\end{equation}
は次数付き自己同型で、$\sigma^2=\iota$、ただし
$\iota(u)=u,\iota(v)=v,\iota(c)=-c,\iota(d)=-d$ である。
命題~\ref{prop:noncubic-relative-serre-on-perfect} の双加群計算から
$S=(-)_\iota(-2)[2]$ であり、
\begin{equation}\label{eq:nchh-good-root}
s=(-)_\sigma(1),\qquad s^2=S^{-1}[2]
\end{equation}
は適合根である。同じ根の AS4 表示は定理~\cite[\texttt{thm:ore-anticommuting-central}]{QuadricIntrinsic2026} の Ore/Zhang 構成が与える。

\begin{proposition}[根の作用と交代 Gale 双対]\label{prop:nchh-root-gale}
\eqref{eq:nchh-good-root} の Hochschild 作用は、前節の Grassmann 模型で
\begin{equation}\label{eq:nchh-involution}
\tau(\mathscr L)=J\mathscr L^{\perp},\qquad
J(X,Y)=(D_0Y,-D_0X),\qquad D_0=\operatorname{diag}(1,-1)
\end{equation}
の微分により与えられる。$X,Y$ を二次正方行列として書いている。
$J:H^{\vee}\to H$ は交代同型であり、中央二重平方の点 $\mathscr L$ は
$\tau$ の固定点、すなわち対応する交代形式の Lagrangian 四平面である。
\end{proposition}

\begin{proof}
行列分解を
\[
\Phi_c=\begin{pmatrix}c&u\\-v&-c\end{pmatrix},\qquad
\Phi_d=\begin{pmatrix}d&\ell_1\\-\ell_2&-d\end{pmatrix}
\]
とする。$M_c=\operatorname{coker}\Phi_c$、$M_d=\operatorname{coker}\Phi_d$ とし、
normal foundation を
\[
(\OO,\pi M_c,\OO(1),\pi M_d(1))
\]
と取る。
$A,C$ はそれぞれ $M_c,M_d$ の二つの次数零生成元の空間、
$B$ は $\Phi_c$ の二つの行、$K$ は $\Phi_d$ の二つの関係列である。
従って \eqref{eq:nchh-delta} の係数行列は
\begin{equation}\label{eq:nchh-plane}
\lambda\longmapsto
\bigl(-\Phi_c^{t}(\lambda),\Phi_d(\lambda)\bigr)
\quad\bigl(\lambda\in\langle u,v,c,d\rangle^{\vee}\bigr).
\end{equation}

二項移動の $D'=H/\mathscr L$ を $\langle u,v,c,d\rangle$ と同定する写像は、
二つの summand 上で次の係数行列を持つ。
\begin{equation}\label{eq:nchh-quotient-identification}
h_0=-D_0\Phi_c=
\begin{pmatrix}-c&-u\\-v&-c\end{pmatrix},\qquad
h_1=\Phi_d^tD_0=
\begin{pmatrix}d&\ell_2\\\ell_1&d\end{pmatrix}.
\end{equation}
実際、$h_0F^{\vee}=h_1j$ であり、両辺は
$u^{\vee}\mapsto-v$, $v^{\vee}\mapsto u$, $c^{\vee},d^{\vee}\mapsto0$
である。ここで $\alpha^2+\beta\gamma=1$ を使う。
$h_0,h_1$ の像は $u,v,c,d$ を張るため、商からの写像は同型である。
$F'$ は $h_1$、$K'$ の二つの関係列は $h_0=-D_0\Phi_c$ に対応する。
これらは移動した実際の Hom と関係の行列である。

右加群の規約も確認する。
$(M_c)_\sigma$ の提示は
$\sigma^{-1}\Phi_c=-D_0\Phi_dD_0$ であり、
生成元を $D_0$ 倍すると $M_d$ へ移る。
一方 $\sigma^{-1}\Phi_d=\Phi_c$ なので $(M_d)_\sigma\simeq M_c$ である。
また canonical 四移動 $t=(-)_\iota(2)$ では
$(M_c)_\iota\simeq M_c$ の生成元比較が $D_0$ となる。
\eqref{eq:nchh-quotient-identification} の $h_1$ は
$K^{\vee}$ を $D_0$ 倍した $\Phi_d$ の行へ対応させ、
$K'$ は $-D_0\Phi_c$ の関係列である。
従って $s$ による四つの二次元空間の比較は、順に
\begin{equation}\label{eq:nchh-four-comparisons}
(D_0,-I_2,D_0,-I_2)
\end{equation}
となる。
この比較を \eqref{eq:nchh-gale} に入れると、第一 summand は $D_0Y$、
第二 summand は $-D_0X$ となり、\eqref{eq:nchh-involution} を得る。

直接の固定点確認もできる。
$\mathscr L^{\perp}$ の元は
\[
X=\begin{pmatrix}a&-\beta r-\alpha s\\\alpha r-\gamma s&a\end{pmatrix},\qquad
Y=\begin{pmatrix}b&r\\s&b\end{pmatrix}
\]
の形である。$J(X,Y)$ は \eqref{eq:nchh-plane} における
$(\lambda_u,\lambda_v,\lambda_c,\lambda_d)=(s,r,-b,-a)$ の像である。
従って $\tau(\mathscr L)=\mathscr L$。
行ごとの座標で $D_0$ の左乗法は対称行列
$D_L=\operatorname{diag}(1,1,-1,-1)$ であり、
\[
J=\begin{pmatrix}0&D_L\\-D_L&0\end{pmatrix},\qquad J^t=-J.
\]
以上は foundation の全 Artin 変形を移動してから、閉ファイバーを指定した
$s$ で戻す構成である。従ってその接作用は、Morita 不変な
Hochschild 作用そのものである。
\end{proof}

\subsection{全 Hochschild 接空間への作用}

\begin{theorem}\label{thm:noncubic-hh-trivial-root}
任意の noncubic $\mathcal H_2$ 対象の有限 foundation を $\Gamma$ とする。
補題~\ref{lem:nchh-exchange} の trace 零の交換対合を用い、
根対が一致する場合には $L_0=D_0$ と選ぶ。
この選択による適合根 $s$ は
\[
s_*|_{\operatorname{HH}^1(\Gamma)}=1,\qquad s_*|_{\operatorname{HH}^2(\Gamma)}=1
\]
を満たす。
二つの根対が異なる場合は $\dim\operatorname{HH}^1=0$, $\dim\operatorname{HH}^2=3$、
一致する場合はそれぞれ $1,4$ である。
\end{theorem}

\begin{proof}
二つの接複体を同変に比較する。
Grassmannian の接空間 $T$ は十六次元である。
\eqref{eq:nchh-involution} は Lagrangian 点における交代 annihilator involution なので
\[
T^+\simeq\operatorname{Sym}^2\mathscr L^{\vee},\qquad
T^-\simeq\bigwedge^2\mathscr L^{\vee},\qquad
(\dim T^+,\dim T^-)=(10,6).
\]
実際、Lagrangian complement を選んで近傍を行列 $Z$ の graph と書けば、
微分は $Z\mapsto Z^t$ である。

次に基底変更群を調べる。
$\operatorname{Lie}\widetilde G=\mathfrak{gl}(A)\oplus\mathfrak{gl}(B)
\oplus\mathfrak{gl}(C)\oplus\mathfrak{gl}(K)$ 上の作用は
\begin{equation}\label{eq:nchh-lie-action}
(a,b,c,k)\longmapsto
(D_0cD_0,-k^t,D_0aD_0,-b^t).
\end{equation}
二つの対を交換するので、その正負固有空間の次元はともに八である。
自明な中心は、scalar 四つ組について
\[
-a-b=c-k
\]
という一つの条件を満たす三次元空間である。
その正固有空間は $a=c$, $b=-k$, $a+b=0$ により一次元、
負固有空間は $a=-c$, $b=k$ により二次元である。
従って有効群の Lie 代数 $\mathfrak g$ について
\begin{equation}\label{eq:nchh-lie-dim}
(\dim\mathfrak g^+,\dim\mathfrak g^-)=(7,6).
\end{equation}

foundation の全変形を表すことから、同変な完全列
\begin{equation}\label{eq:nchh-tangent-sequence}
0\longrightarrow\operatorname{HH}^1(\Gamma)\longrightarrow\mathfrak g
 \longrightarrow T\longrightarrow\operatorname{HH}^2(\Gamma)\longrightarrow0
\end{equation}
がある。これは有向代数の基底変更を除いた結合代数変形の接複体である。
前節の Grassmann 模型は $G$ の全射性と $\delta$ の単射性という開条件だけで得られ、
この完全列に横断方向を除く追加条件は入っていない。
ここで $\OO$ は exceptional であり、
$\operatorname{HH}^0(\Gamma)=Z(\Gamma)\to\operatorname{End}(\OO)$ は $k\to k$ の同型である。
従って pointed 化による追加の一次核はなく、
定理~\ref{thm:artin-foundation-lift} の同定の下で
pointed 慣性の Lie 代数は $\operatorname{HH}^1(\Gamma)$ である。

根対が異なる場合、命題~\ref{prop:noncubic-presentation-stack} の慣性計算により
pointed 慣性は有限である。標数零なので $\operatorname{HH}^1(\Gamma)=0$ である。
\eqref{eq:nchh-tangent-sequence} の負部分は六次元空間から六次元空間への単射であり同型、
従って $\operatorname{HH}^2(\Gamma)^-=0$ となる。正部分は $10-7=3$ 次元である。

根対が一致する場合は $f_1=uv$, $f_2=-uv$, $L_0=D_0$ と取る。
pointed 慣性の連結成分は
$u\mapsto zu$, $v\mapsto z^{-1}v$, $c,d\mapsto c,d$ のトーラスである。
これは $\sigma$ と可換であり、次数移動とも可換なので $s$ の共役作用は恒等である。
従って一次元の $\operatorname{HH}^1(\Gamma)$ は正固有空間に入る。
線型模型でも、$T_z=\operatorname{diag}(z,1)$ に対する安定化群の表示は
\[
(g_A,g_B,g_C,g_K)=(T_z,T_z^{-1},T_z,T_z),
\]
であり、\eqref{eq:nchh-lie-action} で固定される。
負部分は再び同型、正部分の余核は $10-(7-1)=4$ 次元である。
これで両次数への恒等作用が従う。
\end{proof}

\begin{corollary}[全 Artin 変形への根の持上げ]\label{cor:noncubic-root-all-artin}
前定理の根 $s$ は、閉ファイバーを固定した任意の局所 Artin Abel 変形へ、
相対 Serre 平方根として持ち上がる。
その Morita 自己同値の同型類は一意である。
この一意性は平方同型自体の一意な選択を意味しない。
\end{corollary}
\begin{proof}
既存の 定理~\ref{thm:artin-foundation-lift} により、全 Abel 変形を有限自由相対 smooth な
foundation 変形で計算できる。
系~\ref{cor:hhtrivial-root-artin} と定理~\ref{thm:noncubic-hh-trivial-root} の $\operatorname{HH}^2$ 作用から
根の存在を得る。
命題~\ref{prop:artin-serre-root-uniqueness} と $\operatorname{HH}^1$ 作用から同型類の一意性を得る。
補題~\ref{lem:artin-morita-heart-exact} により、持ち上がった Morita 関手は
元の Abel 変形の heart 上で完全である。
本証明は AS4 被覆の横断変形への持上げを与えていない。
\end{proof}

\subsection{全 foundation 変形の Lagrangian 正規形}
\label{subsec:noncubic-lagrangian-slice}

前の接計算は、有限型の具体的な局所表示に強められる。
この表示は foundation の全横断方向を含むが、
それだけで相対 $\mathcal H_2$ の認識を与えるものではない。

四つの二次元空間をそれぞれ $k^2$ と同定し、
第\ref{sec:noncubic-hh-roots}節の $H$ と $J$ を使う。
$\omega$ を行列 $J^{-1}$ の非退化交代形式とする。
$X=\Gr(4,H)$、$L=\operatorname{LG}(4,H,\omega)$ と置けば、
$L$ は $\tau(\mathscr L)=J\mathscr L^\perp$ の固定部分スキームである。
基底変更群の作用は
\[
 \rho(g_A,g_B,g_C,g_K)
 =\operatorname{diag}(g_A^{-t}\otimes g_B^{-t},
                      g_C\otimes g_K^{-t})
\]
である。$X$ 上で自明に作用する中心を
\[
 C_0=\{(aI,bI,cI,dI):a^{-1}b^{-1}=cd^{-1}\}
       \simeq\mathbf G_m^3,
 \qquad G=\widetilde G/C_0
\]
と書く。

\begin{proposition}[Lagrangian な有限型 \texorpdfstring{\'etale}{etale} 表示]
\label{prop:noncubic-lagrangian-etale-slice}
群対合
\begin{equation}\label{eq:noncubic-group-involution}
 \widetilde\theta(g_A,g_B,g_C,g_K)
 =(D_0g_CD_0,g_K^{-t},D_0g_AD_0,g_B^{-t})
\end{equation}
は $G$ の対合 $\theta$ を誘導し、
$\tau(g\mathscr L)=\theta(g)\tau(\mathscr L)$ を満たす。
その固定群の単位成分を $G^+=(G^\theta)^\circ$ と書くと、
\[
 G^+\simeq
 (\GL_2\times\GL_2)/(\mathbf G_m\times\mu_2),
 \qquad \dim G^+=7.
\]
ここで中央部分群の埋込みは
$(z,\epsilon)\mapsto(zI,\epsilon z^{-1}I)$ である。

$G$ の Lie 代数の負固有空間を $\mathfrak g^-$ とする。
自然な軌道微分から得る束写像
\begin{equation}\label{eq:noncubic-normal-orbit-map}
 q:\mathfrak g^-\otimes\mathcal O_L
   \longrightarrow N_{L/X}
   \simeq\bigwedge^2\mathscr U^\vee
\end{equation}
を考える。$\mathscr U$ は $L$ 上の tautological 四平面束である。
その行列式の非消滅開を $L^{\mathrm{tr}}$ とすると、
\[
 [L^{\mathrm{tr}}/G^+]\longrightarrow[X/G]
 \tag{LG-slice}\label{eq:noncubic-lagrangian-stack-map}
\]
は表現可能な \texorpdfstring{\'etale}{etale} 射である。
双方のスタックの次元は三である。
$\det q$ は、定数一次元因子を除けば
Pl\"ucker 線束の三乗 $\mathcal O_L(3)$ の切断である。

中央二重平方の trace 零正規形
\eqref{eq:nchh-plane} が $L^{\mathrm{tr}}$ に属する必要十分条件は
$\alpha\ne0$ である。
根対が異なればこの条件は自動的に成立し、
根対が一致するときは補題~\ref{lem:nchh-exchange} の通り
$L_0=D_0$ と選べば成立する。
従ってその近傍の全 foundation 変形は、
\texorpdfstring{\'etale}{etale} 局所に Lagrangian 四平面で表示できる。
同じ結論は、補題~\ref{lem:nchh-gale} の階数条件を課した
foundation の開部分へ制限したときにも成立する。
\end{proposition}

\begin{proof}
行列を掛けると
\[
 \rho(\widetilde\theta(g))
       =J\rho(g)^{-t}J^{-1}
\]
を得る。この式は $\widetilde\theta^2=1$、
$\widetilde\theta(C_0)=C_0$ と上の同変性を示す。

$\widetilde G$ の固定群は
\[
 (a,b)\longmapsto(a,b,D_0aD_0,b^{-t})
 \qquad(a,b\in\GL_2)
\]
の像である。その $C_0$ との共通部分は
$a=zI$、$b=\epsilon z^{-1}I$、$\epsilon^2=1$ で尽くされる。
従って $G$ 内の像は上で表示した滑らかな連結七次元群である。
標数零の群の固定部分スキームは滑らかであり、その Lie 代数は
$\mathfrak g^+$ である。
\eqref{eq:nchh-lie-dim} によりこの次元は七なので、
上の像は $(G^\theta)^\circ$ と一致する。

$\tau$ の $L$ に沿う微分には、
接束 $T_L\simeq\operatorname{Sym}^2\mathscr U^\vee$ と
法束 $N_{L/X}\simeq\bigwedge^2\mathscr U^\vee$ がそれぞれ
正・負固有束として現れる。
同変性により軌道微分は正負分解を保つため、
\eqref{eq:noncubic-normal-orbit-map} が得られる。
両辺の階数は六であり、その可逆性は $G^+$-不変な開条件である。
$\det(\bigwedge^2\mathscr U^\vee)
=(\det\mathscr U^\vee)^3$ だから、行列式の線束も主張の通りである。

閉部分群 $G^+\subset G$ による associated bundle
$G\times^{G^+}L^{\mathrm{tr}}$ と軌道写像
\[
 m:G\times^{G^+}L^{\mathrm{tr}}\longrightarrow X,
 \qquad[g,\mathscr L]\longmapsto g\mathscr L
\]
を考える。これは滑らかな十六次元空間の間の射である。
$[1,\mathscr L]$ における接写像は
\[
 (\mathfrak g\oplus T_{\mathscr L}L)\big/
 \{(a,-a\mathscr L):a\in\mathfrak g^+\}
 \longrightarrow T_{\mathscr L}X,
 \qquad[(a,v)]\longmapsto a\mathscr L+v.
\]
正部分では恒等、負部分では $q_{\mathscr L}$ となる。
従って $L^{\mathrm{tr}}$ 上では同型である。
$G$-同変性により全ての点で同型なので、$m$ は \texorpdfstring{\'etale}{etale} である。
この射は \eqref{eq:noncubic-lagrangian-stack-map} を atlas $X\to[X/G]$
で基底変更したものだから、同スタック射も表現可能かつ
\texorpdfstring{\'etale}{etale} である。
次元は $10-7=16-13=3$ である。

最後の主張は一般係数で直接確認できる。
$\mathfrak g^-$ の元を
$(a,b,-D_0aD_0,b^t)$、$a,b\in\mathfrak{sl}_2$ で代表し、
各 $\mathfrak{sl}_2$ の基底を
$\operatorname{diag}(1,-1),E_{12},E_{21}$ の順に取る。
\eqref{eq:nchh-plane} の係数行列を $M$ とすると、
法方向の交代形式は $M^tJ^{-1}\rho_*(a,b,-D_0aD_0,b^t)M$ である。
$u^\vee,v^\vee,c^\vee,d^\vee$ の順に対する六つの上三角成分を並べれば、
$q_{\mathscr L}$ の行列は
\[
 Q_{\alpha,\beta,\gamma}=
 \begin{pmatrix}
 -2\alpha&0&0&2\alpha&0&0\\
 0&\gamma&-\alpha&0&\alpha&-\gamma\\
 0&0&-1&0&-1&0\\
 0&-\alpha&-\beta&0&\beta&\alpha\\
 0&1&0&0&0&1\\
 -2&0&0&-2&0&0
 \end{pmatrix},
 \qquad
 \det Q_{\alpha,\beta,\gamma}
       =-32\alpha(\alpha^2+\beta\gamma)=-32\alpha.
\]
一方 $\alpha=0$ なら $\beta\gamma=1$ なので
$f_2=(\beta v)(\gamma u)=uv=f_1$ である。
従って根対非一致の全場合には $\alpha\ne0$ であり、
一致の場合にも対角選択 $\alpha=1,\beta=\gamma=0$ が使える。
foundation を再構成するための追加条件は有限写像の階数条件で開であり、
それらを満たす点の近傍に上の射を制限すればよい。
\end{proof}

\begin{remark}[この表示が与えるもの]
\label{rem:noncubic-lagrangian-scope}
上の射の像は foundation スタックの開部分であり、
全 noncubic 点を含む。その smooth 性は有限代数表示についての主張である。
実際の Morita 平方根との比較には
系~\ref{cor:noncubic-etale-morita-root} を使い、
必要に応じて二つの \texorpdfstring{\'etale}{etale} 近傍の積を取る。
全正軌道 Hom、相対 heart、同じ根の AS4 被覆の存在を、
この Lagrangian 正規形だけから結論してはならない。
特に \eqref{eq:noncubic-lagrangian-stack-map} の標的を
内在的 $\mathfrak H_2$ に置き換えたとは主張していない。
\end{remark}

\begin{remark}[trace 零だけでは良い根を特徴付けない]
\label{rem:noncubic-offdiagonal-root}
根対一致の場合、$L_0=\left(\begin{smallmatrix}0&1\\1&0\end{smallmatrix}\right)$
も trace 零の交換対合であるが、上の \texorpdfstring{\'etale}{etale} 開には属さない。
対応する根は連結慣性
$u\mapsto zu,v\mapsto z^{-1}v$ を逆元へ移すため、
$\operatorname{HH}^1$ に $-1$ で作用する。
\eqref{eq:nchh-tangent-sequence} の正負分解より、この場合
\[
 (\dim\operatorname{HH}^2_+,\dim\operatorname{HH}^2_-)=(3,1).
\]
従ってこの根は全ての一次 Abel 変形に延長するわけではない。
本稿で全 Artin 延長に使う良い根は、根対一致時には
$L_0=D_0$ を選ぶという条件を含む。
\end{remark}

\subsection{Hochschild 二次に自明な根の étale 代数化}
\label{subsec:hh2-etale-morita-root}

次の命題は noncubic 点に限らない。
根の存在する空間から étale な切片を取ることにより、
Hochschild 一次への作用を仮定せずに根を代数化する。

\begin{proposition}[Hochschild 二次に自明な根の étale 代数化]
\label{prop:hh2-etale-morita-root}
$S$ を標数零の代数閉体 $k$ 上の局所有限型 scheme とし、
$\Lambda$ を $S$ 上の有限局所自由代数とする。
次を仮定する。
\begin{enumerate}
\item $\operatorname{Perf}(\Lambda)$ は相対 smooth proper であり、
任意の基底変更後も
$\operatorname{HH}^{<0}(\Lambda)=0$、
$\operatorname{HH}^0(\Lambda)=\OO_S$ である。
\item 相対 inverse Serre $t=S_\Lambda^{-1}[2]$ は
可逆 perfect 双加群で表され、任意の導来基底変更と可換する。
\item $x\in S(k)$ と
$p_0\in\operatorname{DPic}_k(\Lambda_x)$ が与えられ、
$p_0^2=t_x$、
$(p_0)_*|_{\operatorname{HH}^2(\Lambda_x)}=1$ である。
\end{enumerate}
$\mathcal G\to S$ を Morita 自己同値の étale 局所同型類の代数空間、
$\tau=[t]$ とし、
\[
 \mathcal Z=
 \mathcal G\times_{\mathrm{sq},\mathcal G,\tau}S
\]
を平方根空間とする。
このとき $\mathcal Z\to S$ は $p_0$ で smooth である。
さらに有限型 affine な étale 近傍 $(U,u)\to(S,x)$ 上に、
可逆 perfect $\Lambda_U$-双加群 $P$ と実際の同型
\[
 P\otimes_{\Lambda_U}^{\mathbf L}P\simeq t_U
 \tag{N-root}\label{eq:hh2-etale-actual-root}
\]
が存在し、$[P_u]=p_0$ となる。
追加で $(p_0)_*|_{\operatorname{HH}^1(\Lambda_x)}=1$ なら、
$\mathcal Z\to S$ 自体が $p_0$ で étale であり、
同じ閉根を持つ Artin 持上げの Morita 同型類は一意である。
\end{proposition}

\begin{proof}
命題~\ref{prop:quadric-etale-morita-root} の代数空間の構成は、
可換 quadric という仮定を使わず、仮定 (1), (2) に適用できる。
詳しくは \cite[著者公開版 April 2007, Corollary 3.24,
Lemma 3.25]{TV2007} により $\mathcal G$ は $S$ 上局所有限表示の代数空間である。
同論文 Corollary 3.22 により、実際の可逆 perfect 双加群のスタックから
$\mathcal G$ への射は étale 局所に中立な $\mathbf G_m$-gerbe である。
$S$ は局所 Noetherian なので $\mathrm{sq}$ も局所有限表示であり、
$\mathcal Z$ は $S$ 上局所有限表示となる。
ここでは基底の線束による tensor が消える étale 層化を用いる。

$p_0$ を閉点とする任意の局所 Artin 変形
$\operatorname{Spec} A\to\mathcal Z$ と、これに整合する小拡大
$A'\twoheadrightarrow A$ 上の $\operatorname{Spec} A'\to S$ を取る。
剰余体は $k$ とする。
$A,A'$ は strictly henselian で、上の gerbe は中立であり、
線束も自明だから、指定された点は実際の
$p_A\in\operatorname{DPic}_A(\Lambda_A)$、
$p_A^2=t_A$ で表せる。
補題~\ref{lem:root-obstruction-composition} により、
$p_A$ の持上げ障害
$b\in\operatorname{HH}^2(\Lambda_x)\otimes_k\ker(A'\to A)$ は
\[
 b+(p_0)_*b=0
\]
を満たす。
仮定 (3) と標数零により $b=0$ である。
さらに命題~\ref{prop:artin-serre-square-correction} は、
指定された $p_A$ を保ったまま平方関係を補正する。
従ってこの小拡大に対し、指定された点を保つ
$\operatorname{Spec} A'\to\mathcal Z$ が存在する。

$(V,v)\to(\mathcal Z,p_0)$ を étale scheme chart とする。
形式的 étale 性により上の持上げは、指定された
$\operatorname{Spec} A\to V$ も保って $V$ へ持ち上がる。
閉点での小拡大判定
\cite[Tag 02HX, Lemma 37.12.1]{StacksSmoothAtlas}
により $V\to S$ は $v$ で smooth であり、
従って $\mathcal Z\to S$ は $p_0$ で smooth である。

$(V,v)$ を smooth locus に縮める。
$v$ は fiber の閉点であり、$\kappa(v)=\kappa(x)=k$ だから、
\cite[Tag 057G, Lemma 37.38.5]{StacksSmoothAtlas}
により、$v$ を含む immersion $U_0\to V$ で
$U_0\to S$ が étale となるものが存在する。
これは指定された根を通る étale な切片であり、
根の一意性を要求しない。
$U_0$ を通じて上の gerbe を étale 局所に中立化し、
実際の可逆双加群 $P$ を取る。
定理~\ref{thm:bimodule-line-comparison} は比較線束
\[
 L=H^0\mathbf R\mathcal Hom_{\Lambda_U^e}
       (P\otimes_{\Lambda_U}^{\mathbf L}P,t_U)
\]
と評価同型
$(P\otimes_{\Lambda_U}^{\mathbf L}P)\otimes L\simeq t_U$
を与える。
$u$ の Zariski 近傍で $L$ を自明化すれば
\eqref{eq:hh2-etale-actual-root} を得る。
各細分は étale であり、$u$ を含む affine 開へ縮めて
$U$ を有限型にできる。
その像は $x$ を含む開部分なので、必要なら $S$ をこの像へ縮める。

追加の Hochschild 一次の仮定の下では、
命題~\ref{prop:artin-serre-root-uniqueness} により
指定された閉根を持つ Artin 持上げの同型類は一意である。
従って smooth な $\mathcal Z\to S$ の $p_0$ での相対接空間は零であり、
その相対次元は零である。よってこの射自体が $p_0$ で étale となる。
平方同型そのものを一意に選ぶという意味ではない。
\end{proof}

\begin{remark}[Hochschild 一次の役割]
\label{rem:hh2-etale-root-tangent}
同じ閉根を持つ根の一次的な差は
$a\in\operatorname{HH}^1(\Lambda_x)$ で
$a+(p_0)_*a=0$ を満たすものである。
従って平方根空間の相対接空間は
$\ker(1+(p_0)_*)$ である。
Hochschild 一次への恒等作用はこの相対接空間を消し、
根の空間自体を étale にするために使う。
étale 近傍上で一つの根を選ぶ存在定理には不要である。
\end{remark}

\begin{corollary}[全 noncubic 点の有限型 étale Morita 根]
\label{cor:noncubic-etale-morita-root}
$S$ を有限 foundation スタックの有限型 affine chart とし、
$x\in S(k)$ を任意の noncubic $\mathcal H_2$ 対象の
normal foundation に対応する点とする。
定理~\ref{thm:noncubic-hh-trivial-root} の良い根 $s_0$ を選ぶ。
このとき $x$ の有限型 affine étale 近傍 $(U,u)\to(S,x)$ 上に、
$s_0$ を閉 fiber とする可逆 perfect 双加群 $P$ と
実際の平方同型 \eqref{eq:hh2-etale-actual-root} が存在する。
平方根空間自体もこの点で étale である。

四つの普遍 projective 対象 $F_0,\ldots,F_3$ を、
閉 fiber で
$(\OO,\pi M_c,\OO(1),\pi M_d(1))$ に対応させる。
近傍をさらに縮めれば、$P$ による自己同値 $\mathsf s$ は
\[
 \mathsf s F_0\simeq F_2,\qquad
 \mathsf s F_1\simeq F_3
 \tag{N-window}\label{eq:noncubic-etale-root-window}
\]
を満たす。これらの同型は任意の導来基底変更に保たれる。
\end{corollary}

\begin{proof}
補題~\ref{lem:foundation-relative-smooth-proper} と
補題~\ref{lem:foundation-relative-antiserre-perfect} が前命題の
相対 smooth proper 性と Serre の基底変更を与える。
任意の基底変更後に中心が基底環であり負の Hochschild 次数が零であることは、
命題~\ref{prop:quadric-etale-morita-root} の有向代数の中心の計算による。
定理~\ref{thm:noncubic-hh-trivial-root} は両 Hochschild 次数への恒等作用を与えるので、
前命題を適用できる。

最後の対象比較は有限個の perfect 複体の開性による。
$a=0,1$ に対し
$\mathbf R\Hom(\mathsf sF_a,F_{a+2})$ は底上 perfect で、
命題~\ref{prop:nchh-root-gale} の実際の根の比較により
閉 fiber は $k[0]$ である。
従って近傍で次数零の線束となる。
閉 fiber の同型に対応する切断を選び、
その射の cone に perfect Nakayama を適用すると
\eqref{eq:noncubic-etale-root-window} が得られる。
この同型は derived tensor 後も同型のままである。
\end{proof}

\begin{remark}[有限型近傍でまだ結論していないこと]
\label{rem:noncubic-etale-root-scope}
上の $U$ は有限代数と perfect Morita 根の実際の有限型族である。
この構成だけから相対 Abel heart、
全正軌道 Hom の同じ開近傍での次数零集中、
全軌道環の平坦性・Noetherian 性・AZ 豊富性、
同じ根に対する相対 AS4 被覆は従わない。
従って $U\to\mathfrak H_2$ や
noncubic 点を含む smooth Artin 開近傍を、ここでは構成したとは主張しない。
明示的な noncubic 二形式表示族の効果性と、
cubic を横断するこの étale 近傍の効果性も区別する。
\end{remark}

\subsection{根の有限性、計算の検証と残る問題}

根対が異なる場合、命題~\cite[\texttt{prop:first-orbit-object}]{QuadricIntrinsic2026} により任意の適合根の
$\OO$ 上の像は同じ内在的対象である。
従って二つの根の差は pointed 自己同値を与え、全根の同型類は
有限な pointed 自己同値群の torsor の部分集合に入る。
一般点でのこの群は既存原稿では $D_8\times C_2$（位数十六）であり、
二つの平方根変数の符号群 $\mu_2^2$ だけに限定してはいけない。
本節は全ての根の作用を恒等と主張せず、上記の trace 零の選択による
一つの根を特定して証明した。
他の根に対する作用は、それとの差である pointed 自己同値の作用との合成である。

補足の厳密有理計算
\texttt{computations/verify\_noncubic\_gale\_root.py} は、非交差、一根共有、一致の各配置で
$\mathscr L$、$J$、Grassmann tangent の十六次元行列、
基底変更軌道の全十六生成元を構成する。
浮動小数点や乱数を使わず、
\[
\begin{array}{c|ccc}
 &\text{非交差}&\text{一根共有}&\text{一致}\\\hline
\dim(\text{orbit}^+,\text{orbit}^-)&(7,6)&(7,6)&(6,6)\\
\dim(\operatorname{HH}^2_+,\operatorname{HH}^2_-)&(3,0)&(3,0)&(4,0)
\end{array}
\]
を確認する。一般の配置への証明はこの三例の計算から外挿せず、
定理~\ref{thm:noncubic-hh-trivial-root} の同変接複体と安定化群の議論による。

\paragraph{同じ根の AS4 被覆は別の問題である.}
持ち上げた同じ根の軌道環の $B_1$ は四次元の $D$ であり、
$B_2=\Hom(E_0,E_4)$ と積は同じ有限 foundation の bar 複体から計算できる。
根を適合するように twist する前の中央二重平方環では、二次双対に正規な二次元の
pencil が存在するが、固定した適合根の $\mathcal B^\sigma$ でも
同じ pencil が残るとは限らない。ambient の関係線の $\sigma$ 不変性が必要である。
非交差・一根共有・一致の三つの有理配置では、三次平坦性を保つ
関係空間の接次元 $18,18,19$ の全方向に対し、選んだ projective 正規線と
その normalizer が一次に一意に持ち上がることを
\texttt{computations/verify\_noncubic\_cover\_tangent.py} の有限線型代数で確認した。
これはこの特定の被覆の一次計算であり、それ自体から高次存在は従わない。
次節では別の中央 AS4 被覆を全配置で構成し、
系~\ref{cor:noncubic-all-artin-h2} で同じ良い根に対する全 Artin 延長を証明する。
有限型近傍の一様な効果性はなお別問題である。
